\subsection{Errors --- MCQs, Problems \& Traps}

\begin{mcqbox}
\textbf{Q1.} An iteration reports $|\epsilon_a| = 0.5\%$. The number of significant digits guaranteed correct is\\
(a) 1 \quad (b) 2 \quad (c) 3 \quad (d) 4\\
\textbf{Ans: (b)} --- $m=\lfloor 2-\log_{10}(2\times0.5)\rfloor = \lfloor 2-0\rfloor = 2$. Memorize the ladder $5\%,0.5\%,0.05\%,0.005\% \to 1,2,3,4$ digits.

\textbf{Q2.} $0.00640$ has how many decimal places and significant figures respectively?\\
(a) 5 and 5 \quad (b) 3 and 3 \quad (c) 5 and 3 \quad (d) 3 and 5\\
\textbf{Ans: (c)} --- decimal places count everything after the point; significant figures start at the first non-zero digit (6, 4, 0).

\textbf{Q3.} The central-difference derivative approximation has truncation error of order\\
(a) $O(h)$ \quad (b) $O(h^2)$ \quad (c) $O(h^3)$ \quad (d) $O(1)$\\
\textbf{Ans: (b)} --- the odd-order Taylor terms cancel. Forward difference is only $O(h)$.

\textbf{Q4.} A forward-difference derivative of $\sin x$ is computed with $h = 10^{-1}, 10^{-2},\dots,10^{-16}$. The error is smallest at roughly\\
(a) $h=10^{-16}$ \quad (b) $h=10^{-8}$ \quad (c) $h=10^{-2}$ \quad (d) it decreases all the way\\
\textbf{Ans: (b)} --- total error is U-shaped with minimum near $h^\ast\sim\sqrt{\varepsilon_{\text{mach}}}\approx10^{-8}$. Beyond that, round-off in the subtraction takes over and the answer gets \emph{worse}.

\textbf{Q5.} Two numbers each accurate to 6 significant digits are subtracted and are nearly equal. What happens?\\
(a) absolute error grows \quad (b) relative error grows sharply \\
(c) both errors shrink \quad (d) nothing changes\\
\textbf{Ans: (b)} --- the absolute error is unchanged but the result is tiny, so the \emph{relative} error explodes. This is subtractive cancellation.

\textbf{Q6.} Which is computable \emph{without} knowing the exact answer?\\
(a) $E_t$ \quad (b) $\epsilon_t$ \quad (c) $\epsilon_a$ \quad (d) all three\\
\textbf{Ans: (c)} --- $\epsilon_a$ compares the present iterate with the previous one, which is why it is the quantity actually monitored.

\textbf{Q7.} In the trunnion case study the constant-$\alpha$ model predicted a contraction of $0.01504''$ against $0.015''$ required, yet the part jammed. This is an example of\\
(a) round-off error \quad (b) truncation error \quad (c) modelling error \quad (d) a programming bug\\
\textbf{Ans: (c)} --- the arithmetic was correct; the assumption that $\alpha$ is constant was not. No amount of extra precision fixes a wrong model.

\textbf{Q8.} Which statement is \textbf{FALSE}?\\
(a) $\epsilon_a$ is undefined at the first iteration \quad (b) $E_t$ can be negative\\
(c) Round-off error can be eliminated by using more iterations \quad (d) Truncation error appears in numerical integration\\
\textbf{Ans: (c)} --- round-off comes from finite precision, not from stopping early; more iterations can actually make it accumulate.
\end{mcqbox}

\begin{probbox}
\PQ{P1.} A computation produces successive estimates $2.0122$ then $2.0136$. (i) Compute $|\epsilon_a|$. (ii) How many significant digits are guaranteed? (iii) What $|\epsilon_a|$ would be needed for 4 digits?

\ans (i) $|\epsilon_a| = \left|\dfrac{2.0136-2.0122}{2.0136}\right|\times100\% = 0.069527\%$.

(ii) $m = \lfloor 2-\log_{10}(2\times0.069527)\rfloor = \lfloor 2-\log_{10}(0.139054)\rfloor = \lfloor 2+0.8567\rfloor = \lfloor 2.8567\rfloor = \mathbf{2}$.
Check against the criterion directly: $0.5\times10^{2-2}\% = 0.5\%$ and $0.069527 \le 0.5$ $\checkmark$; for $m=3$ we would need $\le 0.05\%$, which fails.

(iii) For $m=4$: $|\epsilon_a| \le 0.5\times10^{2-4}\% = \mathbf{0.005\%}$.

\PQ{P2.} Show that the relative-error criterion $\left|\frac{x-\hat x}{x}\right| \le \tfrac12\times10^{-M}$ implies the absolute window $|x-\hat x| \le \tfrac12\times10^{\,n-M+1}$, where $|x| = m\times10^n$ with $1\le m<10$.

\ans Multiply the hypothesis through by $|x| = m\,10^n$:
\[
|x-\hat x| \;=\; \left|\frac{x-\hat x}{x}\right|\cdot m\,10^{n} \;\le\; \tfrac12\times10^{-M}\cdot m\,10^{n} \;=\; \frac{m}{2}\times10^{\,n-M}.
\]
Since normalization forces $m<10$,
\[
|x-\hat x| \;<\; \frac{10}{2}\times10^{\,n-M} \;=\; 5\times10^{\,n-M} \;=\; \tfrac12\times10^{\,n-M+1},
\]
which is exactly the absolute window --- half a unit in the place of the $M$-th significant digit. $\blacksquare$

\PQ{P3.} $\pi = 3.14159265\dots$ is approximated by $\hat x = 3.14$. Verify that this is correct to 3 significant figures, and find the full band of values that would also qualify.

\ans Here $|x| = 3.14159\ldots = m\times10^n$ with $m=3.14159$, $n=0$. For $M=3$ the window is
\[
|x-\hat x| \le \tfrac12\times10^{\,0-3+1} = \tfrac12\times10^{-2} = 0.005 .
\]
Actual error: $|3.14159265-3.14| = 0.00159265 \le 0.005$ $\checkmark$, so $3.14$ is correct to 3 significant figures.
Band: $3.14159265 \pm 0.005 = [\mathbf{3.13659},\ \mathbf{3.14659}]$.
\end{probbox}

\begin{trapbox}
\begin{itemize}
\item \textbf{The dash in the first row.} Any iteration table you are asked to complete has $|\epsilon_a|$ undefined on iteration 1. Filling in a number there is an instant mark lost.
\item \textbf{$0.5\times10^{2-m}\%$ already carries the percent sign.} Converting to a fraction \emph{and} keeping the \% divides by 100 twice.
\item \textbf{``Accurate to 5 decimal places'' $\ne$ ``accurate to 5 significant figures.''} For a number like $0.00640$ these differ by two orders of magnitude.
\item \textbf{Smaller $h$ is not always better.} Below $h^\ast\sim\sqrt{\varepsilon_{\text{mach}}}$ the forward-difference error rises again. An option saying ``error $\to 0$ as $h\to0$'' is a distractor on a real machine.
\item \textbf{$E_t$ vs $E_a$.} If a question supplies only two successive iterates, it cannot be asking for true error --- there is no true value in play.
\item \textbf{Sign of the error.} $E_t = \text{true} - \text{approx}$, so a negative $E_t$ means you \emph{overestimated}. Options often differ only in sign.
\end{itemize}
\end{trapbox}

\subsection{Root Finding --- MCQs, Problems \& Traps}

\begin{mcqbox}
\textbf{Q9.} Bisection is applied on $[2,5]$. The minimum number of iterations guaranteeing an error below $10^{-4}$ is\\
(a) 13 \quad (b) 14 \quad (c) 15 \quad (d) 18\\
\textbf{Ans: (c)} --- $n \ge \log_2(3/10^{-4}) = 14.87 \Rightarrow 15$. (Resulting width $3/2^{15} = 9.16\times10^{-5}$.)

\textbf{Q10.} For bisection, the number of iterations needed to reach a given tolerance depends on\\
(a) the function only \quad (b) the initial bracket and the tolerance only\\
(c) the function and the bracket \quad (d) the derivative of the function\\
\textbf{Ans: (b)} --- $n \ge \log_2\frac{b-a}{\varepsilon}$ contains no reference to $f$ whatsoever.

\textbf{Q11.} Starting from $[1,2]$, the bisection bracket width after 5 iterations is\\
(a) $0.0625$ \quad (b) $0.03125$ \quad (c) $0.1$ \quad (d) $0.015625$\\
\textbf{Ans: (b)} --- $1/2^5 = 0.03125$.

\textbf{Q12.} $f(x) = 1/x$ on $[-2,3]$ satisfies $f(x_\ell)f(x_u)<0$. Bisection will\\
(a) find the root $x=0$ \quad (b) converge to the singularity at $x=0$, which is not a root\\
(c) fail to start \quad (d) diverge\\
\textbf{Ans: (b)} --- the sign change is caused by a discontinuity, not a zero. Bisection converges anyway and reports a wrong answer confidently.

\textbf{Q13.} $f(x)=x^2$ has a root at $x=0$. Bisection\\
(a) converges quadratically \quad (b) needs one guess only\\
(c) cannot be started, since no sign change exists \quad (d) converges in one step\\
\textbf{Ans: (c)} --- the curve touches the axis without crossing, so no bracket with $f(x_\ell)f(x_u)<0$ exists anywhere.

\textbf{Q14.} ``Newton--Raphson converges quadratically'' means\\
(a) the error halves each step \quad (b) the error is squared each step, so correct digits roughly double\\
(c) it needs two starting guesses \quad (d) it converges in two iterations\\
\textbf{Ans: (b)} --- $e_{k+1}\approx\frac{f''}{2f'}e_k^2$.

\textbf{Q15.} Newton--Raphson applied to $f(x)=x^2$ from $x_0=1$ produces\\
(a) $0.5, 0.25, 0.125,\dots$ (linear) \quad (b) $0.5, 0.0625,\dots$ (quadratic)\\
(c) $0$ immediately \quad (d) divergence\\
\textbf{Ans: (a)} --- the step simplifies to $x_{k+1} = x_k - \frac{x_k^2}{2x_k} = \frac{x_k}{2}$. A double root destroys quadratic convergence.

\textbf{Q16.} To restore quadratic convergence at a root of multiplicity $k$, the iteration is modified to\\
(a) $x - \frac{f}{k f'}$ \quad (b) $x - k\frac{f}{f'}$ \quad (c) $x - \frac{f'}{f}$ \quad (d) $x - \frac{f}{f''}$\\
\textbf{Ans: (b)}

\textbf{Q17.} Newton's method applied to $f(x) = x^2 - a$ gives the iteration\\
(a) $x_{k+1} = \frac{a}{x_k}$ \quad (b) $x_{k+1} = \frac12\left(x_k + \frac{a}{x_k}\right)$ \quad (c) $x_{k+1} = x_k(2-ax_k)$ \quad (d) $x_{k+1} = 2x_k - a$\\
\textbf{Ans: (b)} --- the classical square-root (``Babylonian'') iteration.

\textbf{Q18.} Which statement about false position is \textbf{FALSE}?\\
(a) It is a bracketing method \quad (b) The bracket width always tends to zero\\
(c) One endpoint may remain fixed for many iterations \quad (d) It uses a straight line through the two bracket endpoints\\
\textbf{Ans: (b)} --- with a frozen endpoint the bracket width stalls, which is why $|\epsilon_a|$ rather than bracket width must be the stopping test.

\textbf{Q19.} The Intermediate Value Theorem, applied with $f(x_\ell)f(x_u)<0$, guarantees\\
(a) exactly one root \quad (b) at least one root \quad (c) at most one root \quad (d) an even number of roots\\
\textbf{Ans: (b)}

\textbf{Q20.} In the floating-ball problem $f'(x) = 3x(x-0.11)$. Which initial guess must be avoided for Newton--Raphson?\\
(a) $x_0 = 0.05$ \quad (b) $x_0 = 0.06$ \quad (c) $x_0 = 0.11$ \quad (d) $x_0 = 0.07$\\
\textbf{Ans: (c)} --- $f'(0.11) = 0$, giving a horizontal tangent and division by zero (as does $x_0=0$).
\end{mcqbox}

\begin{probbox}
\PQ{P4.} Apply \textbf{bisection} to $f(x) = x^3-x-2$ on $[1,2]$ for 5 iterations. Tabulate $x_\ell, x_u, x_m, f(x_m)$ and $|\epsilon_a|$. State how many significant digits are certified at iteration 5, and confirm the width against the $L_n$ formula.

\ans $f(1) = -2 < 0$, $f(2) = 4 > 0$, so a root is bracketed.
\begin{center}\small
\begin{tabular}{@{}cccccc@{}}
\toprule
Iter & $x_\ell$ & $x_u$ & $x_m$ & $f(x_m)$ & $|\epsilon_a|\%$\\
\midrule
1 & 1.000000 & 2.000000 & 1.500000 & $-0.125000$ & ---\\
2 & 1.500000 & 2.000000 & 1.750000 & $+1.609375$ & 14.2857\\
3 & 1.500000 & 1.750000 & 1.625000 & $+0.666016$ & 7.6923\\
4 & 1.500000 & 1.625000 & 1.562500 & $+0.252197$ & 4.0000\\
5 & 1.500000 & 1.562500 & 1.531250 & $+0.059113$ & 2.0408\\
\bottomrule
\end{tabular}
\end{center}
Digits at iteration 5: $m = \lfloor 2-\log_{10}(2\times2.0408)\rfloor = \lfloor 2-0.6107\rfloor = \lfloor1.389\rfloor = \mathbf{1}$ --- only one digit after five iterations, which is exactly bisection's reputation.

Width check: $L_5 = (2-1)/2^5 = 0.03125$, and indeed $1.562500-1.531250 = 0.03125$ $\checkmark$. (True root $1.5213797$.)

\PQ{P5.} Apply \textbf{false position} to the same $f(x)=x^3-x-2$ on $[1,2]$ for 4 iterations. Comment on which endpoint moves.

\ans
\begin{center}\small
\begin{tabular}{@{}ccccc@{}}
\toprule
Iter & $x_L$ & $x_U$ & $x_r$ & $f(x_r)$\\
\midrule
1 & 1.000000 & 2.000000 & 1.333333 & $-9.630\times10^{-1}$\\
2 & 1.333333 & 2.000000 & 1.462687 & $-3.333\times10^{-1}$\\
3 & 1.462687 & 2.000000 & 1.504019 & $-1.018\times10^{-1}$\\
4 & 1.504019 & 2.000000 & 1.516331 & $-2.989\times10^{-2}$\\
\bottomrule
\end{tabular}
\end{center}
$|\epsilon_a|$ at iteration 4 $= |(1.516331-1.504019)/1.516331|\times100\% = 0.8119\%$.

\textbf{Observation:} $f(x_r)$ is negative every time, so $x_r$ always replaces $x_L$; the right endpoint stays \textbf{frozen at $x_U = 2$ forever}. The bracket width only falls from $1$ to $0.4837$ in four iterations even though $x_r$ is converging nicely --- the textbook illustration of why bracket width is a useless stopping test here.

\PQ{P6.} Apply \textbf{Newton--Raphson} to $f(x)=x^3-x-2$ from $x_0 = 1.5$. Do 3 iterations, report $|\epsilon_a|$ and the certified digits each step, and comment on the convergence rate.

\ans $f'(x) = 3x^2-1$.
\[
x_1 = 1.5 - \frac{1.5^3-1.5-2}{3(1.5)^2-1} = 1.5 - \frac{-0.125}{5.75} = 1.52173913
\]
\begin{center}\small
\begin{tabular}{@{}cccc@{}}
\toprule
$k$ & $x_k$ & $|\epsilon_a|\%$ & digits $m$\\
\midrule
1 & 1.52173913 & 1.428571 & 1\\
2 & 1.52137981 & 0.023618 & 3\\
3 & 1.52137971 & 0.000007 & 6\\
\bottomrule
\end{tabular}
\end{center}
Certified digits go $1 \to 3 \to 6$: roughly \textbf{doubling each step}, the signature of quadratic convergence. Compare with P4, where bisection bought a single digit in five iterations.

\PQ{P7.} Prove that Newton--Raphson on $f(x)=x^2$ (double root at $0$) converges only linearly, and give the exact rate.

\ans
\[
x_{k+1} = x_k - \frac{f(x_k)}{f'(x_k)} = x_k - \frac{x_k^2}{2x_k} = x_k - \frac{x_k}{2} = \frac{x_k}{2}.
\]
Hence $e_{k+1} = \tfrac12 e_k$ exactly: convergence is \textbf{linear with rate $1/2$} --- identical to bisection, and nowhere near quadratic. From $x_0=1$ the iterates are $0.5, 0.25, 0.125, 0.0625, \dots$

This matches the general rule $e_{k+1}\approx\left(1-\frac1k\right)e_k$ for multiplicity $k$: here $k=2$ gives rate $1-\frac12 = \frac12$ $\checkmark$. The modified step $x - 2\frac{f}{f'} = x_k - x_k = 0$ lands on the root in a single iteration.

\PQ{P8.} Verify Newton's quadratic error law numerically for $f(x) = x^2-2$ from $x_0 = 1$.

\ans The law predicts $e_{k+1}\approx C e_k^2$ with
\[
C = \left|\frac{f''(x^\ast)}{2f'(x^\ast)}\right| = \frac{2}{2\cdot 2\sqrt2} = \frac{1}{2\sqrt2} = 0.353553 .
\]
The iteration is the Babylonian one, $x_{k+1} = \tfrac12(x_k + 2/x_k)$:
\begin{center}\small
\begin{tabular}{@{}ccccc@{}}
\toprule
$k$ & $x_k$ & actual $e_k$ & predicted $C\,e_{k-1}^2$ & match?\\
\midrule
1 & 1.500000000 & $8.579\times10^{-2}$ & $6.066\times10^{-2}$ & rough (far from root)\\
2 & 1.416666667 & $2.453\times10^{-3}$ & $2.602\times10^{-3}$ & good\\
3 & 1.414215686 & $2.124\times10^{-6}$ & $2.128\times10^{-6}$ & excellent\\
4 & 1.414213562 & $1.595\times10^{-12}$ & $1.595\times10^{-12}$ & exact to 4 s.f.\\
\bottomrule
\end{tabular}
\end{center}
The prediction becomes essentially exact once $x_k$ is close enough for the Taylor argument to hold. Note the error exponents: $10^{-2}\to10^{-3}\to10^{-6}\to10^{-12}$ --- \textbf{doubling the correct digits every step}.
\end{probbox}

\begin{trapbox}
\begin{itemize}
\item \textbf{``Which method is fastest?'' has no universal answer.} Newton is fastest \emph{when it converges}; on $\ln x$ over $[10^{-4},10^4]$ false position (126 iterations) lost badly to bisection (34).
\item \textbf{Bisection iteration count ignores $f$.} Any option that involves the function is wrong by construction.
\item \textbf{Sign change is sufficient, not necessary.} $x^2$ has a root with no sign change; $1/x$ has a sign change with no root. Both appear as MCQ stems.
\item \textbf{Quadratic $\ne$ ``two iterations''} and $\ne$ ``error halves.'' It means the error is \emph{squared}.
\item \textbf{A double root silently kills Newton's speed.} If the stem says ``multiplicity 2'' or gives $f = (x-r)^2 g(x)$, the answer is linear, rate $1/2$.
\item \textbf{$f'(x_0)=0$ is fatal, and near-zero is almost as bad.} The floating-ball bracket ends $0$ and $0.11$ both have $f'=0$.
\item \textbf{In false position, do not use bracket width to stop.} A frozen endpoint keeps it large forever; use $|\epsilon_a|$ on successive $x_r$.
\item \textbf{The reported root must be checked against the interval.} Newton can jump entirely out of your region of interest (root jumping on $\sin x$, converging to $0$ instead of $2\pi$).
\item \textbf{Counting iterations vs counting function evaluations.} Newton needs $f$ \emph{and} $f'$ each step; a question about ``work'' may not have the same answer as one about ``iterations.''
\end{itemize}
\end{trapbox}
